% !TEX root = ../proyect.tex

\chapter{Estructura de la base de datos}\label{cap:estructura-bd}

Este capítulo describe el diseño y la estructura de la base de datos de LinceUS Assistant. La base de datos constituye el núcleo del sistema, almacenando toda la información académica necesaria para responder a las consultas de los estudiantes.

\section{Diseño conceptual}\label{sec:diseno-conceptual}

El diseño de la base de datos se ha estructurado siguiendo el modelo relacional, organizando la información en entidades que representan los elementos principales del dominio universitario: universidades, centros, titulaciones, asignaturas, profesores y horarios.

\subsection{Principios de diseño}

El esquema de base de datos se ha diseñado siguiendo los siguientes principios:

\begin{enumerate}
    \item \textbf{Normalización}: Las tablas están normalizadas hasta la tercera forma normal (3NF) para evitar redundancia y mantener la integridad referencial
    \item \textbf{Escalabilidad}: El diseño permite añadir nuevas universidades, centros y titulaciones sin modificar la estructura
    \item \textbf{Flexibilidad}: Campos de metadatos en formato JSONB permiten almacenar información variable sin alterar el esquema
    \item \textbf{Búsqueda semántica}: Tablas de chunks con vectores de embeddings para búsqueda basada en similitud
    \item \textbf{Trazabilidad}: Campos de timestamp (created\_at, updated\_at) en todas las tablas principales
\end{enumerate}

\section{Modelo entidad-relación}\label{sec:modelo-er}

% El diagrama entidad-relación completo del sistema se muestra en la Figura \ref{fig:diagrama-er}.
A continuación se detallan las principales entidades y sus relaciones del sistema.

% TODO: Descomentar cuando se genere la imagen
% \figura{0.95}{img/diagrama_er}{Diagrama entidad-relación de la base de datos LinceUS}{fig:diagrama-er}{}

\section{Entidades principales}\label{sec:entidades}

\subsection{Universidades}

La tabla \texttt{UNIVERSIDADES} almacena la información básica de cada universidad. Aunque el sistema está diseñado inicialmente para la Universidad de Sevilla, la estructura permite incorporar múltiples instituciones.

\cuadro{|l|l|p{6cm}|}{Estructura de la tabla UNIVERSIDADES}{tab:universidades}{
\textbf{Campo} & \textbf{Tipo} & \textbf{Descripción} \\ \hline
id & uuid & Identificador único (PK) \\ \hline
codigo & varchar & Código de la universidad (e.g., "US") \\ \hline
nombre & varchar & Nombre completo de la universidad \\ \hline
nombre\_corto & varchar & Nombre abreviado \\ \hline
dominio\_web & varchar & Dominio web oficial (e.g., "us.es") \\ \hline
ciudad & varchar & Ciudad donde se ubica \\ \hline
activa & boolean & Indica si está operativa \\ \hline
created\_at & timestamptz & Fecha de creación del registro \\ \hline
updated\_at & timestamptz & Fecha de última actualización \\ \hline
}

\subsection{Centros}

Los centros representan las escuelas y facultades de cada universidad (e.g., ETSII, Facultad de Química).

\cuadro{|l|l|p{6cm}|}{Estructura de la tabla CENTROS}{tab:centros}{
\textbf{Campo} & \textbf{Tipo} & \textbf{Descripción} \\ \hline
id & uuid & Identificador único (PK) \\ \hline
universidad\_id & uuid & Referencia a UNIVERSIDADES (FK) \\ \hline
codigo & varchar & Código del centro (e.g., "ETSII") \\ \hline
codigo\_sevius & varchar & Código del centro en Sevius \\ \hline
codigo\_us & varchar & Código oficial del centro en la US \\ \hline
nombre & varchar & Nombre completo del centro \\ \hline
nombre\_corto & varchar & Nombre abreviado \\ \hline
direccion & varchar & Dirección física del centro \\ \hline
telefono & varchar & Teléfono de contacto \\ \hline
email & varchar & Email de contacto \\ \hline
web & varchar & URL del sitio web del centro \\ \hline
activo & boolean & Indica si está operativo \\ \hline
created\_at & timestamptz & Fecha de creación del registro \\ \hline
updated\_at & timestamptz & Fecha de última actualización \\ \hline
}

\subsection{Departamentos}

Los departamentos son unidades organizativas que agrupan profesores e imparten asignaturas.

\cuadro{|l|l|p{6cm}|}{Estructura de la tabla DEPARTAMENTOS}{tab:departamentos}{
\textbf{Campo} & \textbf{Tipo} & \textbf{Descripción} \\ \hline
id & uuid & Identificador único (PK) \\ \hline
universidad\_id & uuid & Referencia a UNIVERSIDADES (FK) \\ \hline
centro\_id & uuid & Referencia a CENTROS (FK) \\ \hline
codigo & varchar & Código del departamento \\ \hline
codigo\_us & varchar & Código oficial del departamento en la US \\ \hline
nombre & varchar & Nombre completo \\ \hline
siglas & varchar & Siglas del departamento (e.g., "LSI") \\ \hline
email & varchar & Email de contacto \\ \hline
web & varchar & URL del sitio web \\ \hline
activo & boolean & Indica si está operativo \\ \hline
created\_at & timestamptz & Fecha de creación del registro \\ \hline
updated\_at & timestamptz & Fecha de última actualización \\ \hline
}

\subsection{Titulaciones}

Las titulaciones representan los grados, másteres y doctorados ofrecidos por los centros.

\cuadro{|l|l|p{6cm}|}{Estructura de la tabla TITULACIONES}{tab:titulaciones}{
\textbf{Campo} & \textbf{Tipo} & \textbf{Descripción} \\ \hline
id & uuid & Identificador único (PK) \\ \hline
centro\_id & uuid & Referencia a CENTROS (FK) \\ \hline
codigo & varchar & Código interno (e.g., "GII-IS") \\ \hline
codigo\_oficial & varchar & Código RUCT oficial \\ \hline
nombre & varchar & Nombre completo de la titulación \\ \hline
nombre\_corto & varchar & Nombre abreviado \\ \hline
tipo & varchar & GRADO, MASTER o DOCTORADO \\ \hline
plan\_estudios\_anio & int & Año del plan de estudios vigente \\ \hline
creditos\_totales & int & Créditos ECTS totales \\ \hline
duracion\_anios & int & Duración en años \\ \hline
requisito\_idioma & varchar & Nivel de idioma requerido (e.g., "B1") \\ \hline
activa & boolean & Indica si está vigente \\ \hline
created\_at & timestamptz & Fecha de creación del registro \\ \hline
updated\_at & timestamptz & Fecha de última actualización \\ \hline
}

\subsection{Asignaturas}

Las asignaturas son las materias que componen cada titulación.

\cuadro{|l|l|p{5.5cm}|}{Estructura de la tabla ASIGNATURAS}{tab:asignaturas}{
\textbf{Campo} & \textbf{Tipo} & \textbf{Descripción} \\ \hline
id & uuid & Identificador único (PK) \\ \hline
titulacion\_id & uuid & Referencia a TITULACIONES (FK) \\ \hline
departamento\_id & uuid & Referencia a DEPARTAMENTOS (FK) \\ \hline
codigo & varchar & Código de la asignatura \\ \hline
nombre & varchar & Nombre oficial de la asignatura \\ \hline
curso & int & Curso en el que se imparte (1-4) \\ \hline
creditos & decimal & Créditos ECTS \\ \hline
duracion & varchar & A (anual), C1 (1er cuatr.), C2 (2º cuatr.) \\ \hline
tipologia & varchar & TRONCAL, OBLIGATORIA, OPTATIVA \\ \hline
es\_formacion\_basica & boolean & Si es de formación básica \\ \hline
es\_optativa & boolean & Si es optativa \\ \hline
nombre\_normalizado & varchar & Nombre sin tildes ni mayúsculas \\ \hline
palabras\_clave & text[] & Array de palabras clave \\ \hline
activa & boolean & Indica si está vigente \\ \hline
created\_at & timestamptz & Fecha de creación del registro \\ \hline
updated\_at & timestamptz & Fecha de última actualización \\ \hline
}

El campo \texttt{nombre\_normalizado} facilita la búsqueda tolerante a errores, mientras que \texttt{palabras\_clave} permite búsquedas por términos relacionados.

\subsection{Profesores}

La tabla de profesores almacena la información del personal docente.

\cuadro{|l|l|p{6cm}|}{Estructura de la tabla PROFESORES}{tab:profesores}{
\textbf{Campo} & \textbf{Tipo} & \textbf{Descripción} \\ \hline
id & uuid & Identificador único (PK) \\ \hline
departamento\_id & uuid & Referencia a DEPARTAMENTOS (FK) \\ \hline
centro\_id & uuid & Referencia a CENTROS (FK) \\ \hline
nombre & varchar & Nombre del profesor \\ \hline
apellidos & varchar & Apellidos del profesor \\ \hline
nombre\_completo & varchar & Nombre completo (generado) \\ \hline
nombre\_normalizado & varchar & Versión normalizada para búsqueda \\ \hline
categoria\_academica & varchar & Catedrático, Titular, Contratado Doctor, etc. \\ \hline
email & varchar & Email de contacto \\ \hline
telefono & varchar & Teléfono de contacto \\ \hline
despacho & varchar & Número de despacho (e.g., "F1.45") \\ \hline
edificio & varchar & Edificio donde se ubica \\ \hline
planta & varchar & Planta del despacho \\ \hline
web\_personal & varchar & URL de página personal \\ \hline
enlace\_perfil & varchar & URL del perfil en la web del departamento \\ \hline
orcid & varchar & Identificador ORCID \\ \hline
activo & boolean & Indica si está en activo \\ \hline
created\_at & timestamptz & Fecha de creación del registro \\ \hline
updated\_at & timestamptz & Fecha de última actualización \\ \hline
}

\subsection{Horarios}

Los horarios relacionan grupos de clase con aulas, profesores y franjas horarias.

\cuadro{|l|l|p{6cm}|}{Estructura de la tabla HORARIOS}{tab:horarios}{
\textbf{Campo} & \textbf{Tipo} & \textbf{Descripción} \\ \hline
id & uuid & Identificador único (PK) \\ \hline
grupo\_id & uuid & Referencia a GRUPOS\_CLASE (FK) \\ \hline
aula\_id & uuid & Referencia a AULAS (FK) \\ \hline
profesor\_id & uuid & Referencia a PROFESORES (FK) \\ \hline
dia\_semana & int & Día de la semana (1-5) \\ \hline
hora\_inicio & time & Hora de inicio de la clase \\ \hline
hora\_fin & time & Hora de fin de la clase \\ \hline
fecha\_inicio & date & Fecha de inicio del periodo \\ \hline
fecha\_fin & date & Fecha de fin del periodo \\ \hline
cuatrimestre & varchar & Cuatrimestre académico (C1, C2 o anual) \\ \hline
notas & text & Observaciones \\ \hline
activo & boolean & Indica si está vigente \\ \hline
created\_at & timestamptz & Fecha de creación del registro \\ \hline
updated\_at & timestamptz & Fecha de última actualización \\ \hline
}

\section{Tablas de vectorización (RAG)}\label{sec:vectorizacion}

Para implementar el sistema RAG (\textit{Retrieval-Augmented Generation}), se han diseñado tablas especializadas que almacenan fragmentos de documentos junto con sus representaciones vectoriales.

\subsection{Planes docentes}

\cuadro{|l|l|p{6cm}|}{Estructura de la tabla PLANES\_DOCENTES}{tab:planes-docentes}{
\textbf{Campo} & \textbf{Tipo} & \textbf{Descripción} \\ \hline
id & uuid & Identificador único (PK) \\ \hline
asignatura\_id & uuid & Referencia a ASIGNATURAS (FK) \\ \hline
curso\_academico & varchar & Curso académico (e.g., "2025-26") \\ \hline
grupo & varchar & Grupo al que corresponde \\ \hline
url\_documento & varchar & URL del PDF original \\ \hline
hash\_documento & varchar & Hash SHA256 del documento \\ \hline
fecha\_documento & date & Fecha del documento \\ \hline
coordinador\_nombre & varchar & Nombre del coordinador \\ \hline
estado\_rag & varchar & pendiente, procesando, completado, error \\ \hline
fecha\_procesamiento & timestamptz & Fecha de procesamiento \\ \hline
error\_procesamiento & text & Mensaje de error si lo hay \\ \hline
created\_at & timestamptz & Fecha de creación del registro \\ \hline
updated\_at & timestamptz & Última actualización \\ \hline
}

\subsection{Chunks de planes docentes}

\cuadro{|l|l|p{5.5cm}|}{Estructura de la tabla PLANES\_DOCENTES\_CHUNKS}{tab:chunks-planes}{
\textbf{Campo} & \textbf{Tipo} & \textbf{Descripción} \\ \hline
id & uuid & Identificador único (PK) \\ \hline
plan\_docente\_id & uuid & Referencia a PLANES\_DOCENTES (FK) \\ \hline
contenido & text & Texto del fragmento \\ \hline
embedding & vector & Vector de 768 dimensiones \\ \hline
seccion & varchar & Sección del documento (e.g., ``Evaluación'') \\ \hline
subseccion & varchar & Subsección si aplica \\ \hline
numero\_pagina & int & Número de página del PDF \\ \hline
orden\_chunk & int & Orden del chunk en el documento \\ \hline
metadata & jsonb & Metadatos adicionales en formato JSON \\ \hline
created\_at & timestamptz & Fecha de creación del registro \\ \hline
updated\_at & timestamptz & Última actualización \\ \hline
}

El campo \texttt{embedding} utiliza el tipo \texttt{vector} proporcionado por la extensión pgvector. Los vectores se generan usando modelos de embeddings especializados y permiten búsquedas de similitud semántica mediante consultas como:

\begin{verbatim}
SELECT contenido, embedding <=> query_vector AS distance
FROM planes_docentes_chunks
ORDER BY distance
LIMIT 5;
\end{verbatim}

\section{Relaciones entre entidades}\label{sec:relaciones}

Las principales relaciones del modelo son:

\begin{itemize}
    \item Una \textbf{universidad} agrupa múltiples \textbf{centros} y \textbf{departamentos} (FK \texttt{universidad\_id} en ambos).
    \item Un \textbf{centro} ofrece múltiples \textbf{titulaciones} (FK \texttt{centro\_id} en \textsc{titulaciones}) y aloja múltiples \textbf{aulas}. Adicionalmente, un departamento queda adscrito a un centro a través de la FK \texttt{centro\_id} en \textsc{departamentos}.
    \item Una \textbf{titulación} contiene múltiples \textbf{asignaturas} (FK \texttt{titulacion\_id}).
    \item Una \textbf{asignatura} es impartida por un \textbf{departamento} (FK \texttt{departamento\_id}) y tiene múltiples \textbf{planes docentes} (FK \texttt{asignatura\_id} en \textsc{planes\_docentes}) y \textbf{grupos de clase}.
    \item Un \textbf{departamento} agrupa a múltiples \textbf{profesores} (FK \texttt{departamento\_id}). El profesor también queda asociado a un \textbf{centro} mediante la FK \texttt{centro\_id} en \textsc{profesores}, lo que facilita las consultas que filtran por escuela sin tener que recorrer el departamento.
    \item Un \textbf{profesor} imparte múltiples \textbf{asignaturas} y está asignado a \textbf{horarios} (FK \texttt{profesor\_id} en \textsc{horarios}).
    \item Un \textbf{grupo de clase} tiene un conjunto de \textbf{horarios} asociados (FK \texttt{grupo\_id}).
    \item Un \textbf{horario} relaciona un \textbf{grupo} (FK \texttt{grupo\_id}), un \textbf{aula} (FK \texttt{aula\_id}) y un \textbf{profesor} (FK \texttt{profesor\_id}), e incluye el \texttt{cuatrimestre} para distinguir las sesiones del mismo grupo en periodos académicos distintos.
    \item Un \textbf{plan docente} se descompone en múltiples \textbf{chunks} para vectorización (FK \texttt{plan\_docente\_id} en \textsc{planes\_docentes\_chunks}, con \texttt{ON DELETE CASCADE}).
\end{itemize}

Todas las tablas principales registran además \texttt{created\_at} y \texttt{updated\_at} (\texttt{timestamptz}) para auditar la trazabilidad temporal de cada registro, en línea con el principio enunciado en la Sección~\ref{sec:diseno-conceptual}.

\section{Índices y optimización}\label{sec:indices}

La estrategia de indexación cubre tres tipos de acceso según el patrón de consulta:

\begin{itemize}
    \item \textbf{B-tree (PKs y FKs)}: Índices únicos en todas las claves primarias (\texttt{\_pkey}) y en las claves foráneas más consultadas (\texttt{titulacion\_id}, \texttt{grupo\_id}, \texttt{profesor\_id}, \texttt{departamento\_id}, etc.), más índices adicionales sobre campos de filtrado frecuente como \texttt{codigo}, \texttt{curso}, \texttt{dia\_semana} o \texttt{estado\_rag}.
    \item \textbf{GIN trigrama}: El campo \texttt{nombre\_normalizado} de \textsc{asignaturas} lleva un índice \texttt{gin\_trgm\_ops} además del B-tree, lo que permite búsquedas \texttt{ILIKE} eficientes sin exploración secuencial. Este índice es la base del \textit{fallback} por palabras clave del pipeline RAG (Sección~\ref{subsec:dd-fallback}).
    \item \textbf{HNSW}: El campo \texttt{embedding} de \textsc{planes\_docentes\_chunks} usa un índice \textit{Hierarchical Navigable Small World}~\cite{malkov2018hnsw} con parámetros $m=16$ y $ef\_construction=64$ (Sección~\ref{subsec:dd-embeddings}), que ofrece búsqueda aproximada de vecinos más cercanos en tiempo sub-lineal.
\end{itemize}

El índice HNSW sobre los chunks de planes docentes se crea así:

\begin{verbatim}
CREATE INDEX idx_chunks_embedding_hnsw
ON planes_docentes_chunks
USING hnsw (embedding vector_cosine_ops)
WITH (m = 16, ef_construction = 64);
\end{verbatim}

\section{Integridad referencial}\label{sec:integridad}

Todas las relaciones entre tablas están protegidas por restricciones de clave foránea con políticas de borrado apropiadas:

\begin{itemize}
    \item \textbf{ON DELETE CASCADE}: En relaciones donde el hijo no tiene sentido sin el padre (e.g., chunks sin plan docente)
    \item \textbf{ON DELETE RESTRICT}: En relaciones donde se debe prevenir el borrado accidental (e.g., no borrar una universidad con centros asociados)
    \item \textbf{ON DELETE SET NULL}: En relaciones opcionales donde el padre puede desaparecer sin eliminar al hijo
\end{itemize}

\section{Estrategia de datos de prueba}\label{sec:datos-prueba}

Para el desarrollo y pruebas del sistema se han creado datos de prueba que incluyen:

\begin{itemize}
    \item Universidad de Sevilla con varios centros (ETSII, Facultad de Informática)
    \item Titulaciones de Ingeniería Informática (IS, IC, TI)
    \item Conjunto representativo de asignaturas de cada titulación
    \item Profesores del Departamento de Lenguajes y Sistemas Informáticos
    \item Horarios del curso 2025-2026
    \item Planes docentes oficiales de las asignaturas vectorizados en el corpus RAG
\end{itemize}

Los datos de prueba permiten validar el funcionamiento del sistema sin necesidad de tener acceso a datos reales completos.
